\section{Lattice Matrix Elements}\label{sec:lattice-details}

The lattice gluon matrix elements of the kaon's ground-state  $|0 (P_z) \rangle$  are calculated for several boost momenta $P_z$ and Wilson-line displacement lengths $z$ using the unpolarized gluon operator first introduced in Ref.~\cite{Balitsky:2019krf},
%
\begin{equation}\label{eq:gluon_operator}
 {\cal O}(z)\equiv\sum_{i\neq z,t}{\cal O}(F^{ti},F^{ti};z)-\sum_{i,j\neq z,t}{\cal O}(F^{ij},F^{ij};z),
\end{equation}
%
where ${\cal O}(F^{\mu\nu}, F^{\alpha\beta};z) = F^\mu_\nu(z)U(z,0)F^\alpha_\beta(0)$, the Wilson link length is denoted by $z$, and the field strength $F_{\mu\nu}$ is given by
%
\begin{equation}\label{field_strength}
 F_{\mu\nu} = \frac{i}{8a^2g_0}\left(\mathcal{P}_{[\mu,\nu]}+\mathcal{P}_{[\nu,-\mu]}+\mathcal{P}_{[-\mu,-\nu]}+\mathcal{P}_{[-\nu,\mu]}\right),
\end{equation}
%
where $a$ is the lattice spacing, $g_0$ is the strong coupling constant, with the plaquette $\mathcal{P_{\mu,\nu}}=U_\mu(x)U_\nu(x+a\hat{\mu})U^\dag_\mu(x+a\hat{\nu})U^\dag_\nu(x)$ and $\mathcal{P_{[\mu,\nu]}}=\mathcal{P}_{\mu,\nu}-\mathcal{P}_{\nu,\mu}$.
The operator $\sum_{i\neq z,t}{\cal O}(F^{ti},F^{zi};z)$ corresponds to the same matching kernel as in Ref.~\cite{Balitsky:2019krf}.
However, this operator is not employed, since it vanishes at $P_z=0$ due to kinematic reasons, which increases the difficulty of obtaining the distributions.
Having calculated the aforementioned matrix elements, we investigate their dependence on Ioffe time, $\nu=zP_z$.

On the lattice, we employ clover valence fermions on three ensembles with $N_f=2+1+1$ highly improved staggered quarks (HISQ)~\cite{Follana:2006rc}, generated by the MILC Collaboration~\cite{MILC:2010pul,Bazavov:2012xda} at two lattice spacings ($a\approx 0.12$ and 0.15~fm) and pion mass of 310~MeV.
The clover quark masses were adjusted to reproduce the lightest strange and light sea pseudoscalar meson utilized by the PNDME collaboration~\cite{Rajan:2017lxk,Bhattacharya:2015wna,Bhattacharya:2015esa,Bhattacharya:2013ehc}.
We apply five steps of HYP smearing~\cite{Hasenfratz:2001hp} for the gluon loops to reduce the statistical noise, as elucidated in Ref.~\cite{Fan:2018dxu}.
In order to reach higher meson boost momenta, we use Gaussian momentum smearing~\cite{Bali:2016lva} for the quark fields.
Table~\ref{table-data} shows the lattice spacing $a$, valence pion mass $M_\pi^\text{val}$ and kaon mass $M_K^\text{val}$, lattice size $L^3\times T$, number of configurations $N_\text{cfg}$, and number of total two-point correlator measurements $N_\text{meas}^\text{2pt}$
for the three ensembles.


\begin{table}[!htbp]
\centering
\begin{tabular}{|c|c|c|}
\hline
  ensemble &  a12m310 & a15m310 \\
\hline
  $a$ (fm)   & $0.1207(11)$  & $0.1510(20)$ \\
\hline
  $M_\pi^\text{val}$ (MeV)  & $311.1(6)$ & $319.1(31)$\\
\hline
  $M_{k}^\text{val}$ (MeV)  & 528.0(5) &  531.7(23)  \\
\hline
  $L^3\times T$   & $24^3\times 64$ & $16^3\times 48$ \\
\hline
  $P_z$ (GeV)   &  $[0,2.14]$ &  $[0,2.05]$ \\
\hline
  $N_\text{cfg}$  & 1013 & 900 \\
\hline
  $N_\text{meas}^\text{2pt}$   & 324,160 & 21,600 \\
\hline
%\hline
\end{tabular}
\caption{
Lattice spacing $a$, valence pion mass $M_\pi^\text{val}$ and kaon mass $M_{k}^\text{val}$, lattice size $L^3\times T$, number of configurations $N_\text{cfg}$, number of total two-point correlator measurements $N_\text{meas}^\text{2pt}$, and separation times $t_\text{sep}$ utilized in three-point correlator fits of $N_f=2+1+1$ clover valence fermions on HISQ ensembles generated by MILC Collaboration and analyzed in this study.
}\label{table-data}
\end{table}




For a meson $\Phi$, the two-point correlator is given by
%
\begin{align}
C_\Phi^\text{2pt}(P_z;t)&=
 \int d^3y\, e^{-i y\cdot P_z} \langle \chi_\Phi(\vec{y},t)|\chi_\Phi(\vec{0},0)\rangle \nonumber \\
  &= |A_{\Phi,0}|^2 e^{-E_{\Phi,0}t} + |A_{\Phi,1}|^2 e^{-E_{\Phi,1}t} + ...,
\label{eq:2pt_fit_formula}
\end{align}
%
where $P_z$ is the meson momentum in the $z$-direction, $\chi_\Phi=\bar{q}_1\gamma_5 q_2$ is the pseudoscalar-meson interpolation operator, and $t$ is the Euclidean time.
$|A_{\Phi,i}|^2$ and $E_{\Phi,i}$ denote the amplitude and energy for the ground state ($i=0$) and the first excited state ($i=1$), respectively.
We generate the three-point gluon correlators by multiplying the kaon two-point correlators with the gluon operators.
The matrix elements of the gluon operators are then extracted by fitting the three-point correlators to the first terms of the energy-eigenstate expansion,
%
\begin{align}\label{eq:3ptC}
&C_\Phi^\text{3pt}(z,P_z;t_\text{sep},t) \nonumber \\
&=\int d^3y\, e^{-iy\cdot P_z}\langle \chi_\Phi(\vec y,t_\text{sep})|{\cal O}(z,t)|\chi_\Phi(\vec 0,0)\rangle \nonumber \\
&= |A_{\Phi,0}|^2\langle 0|{\cal O}|0\rangle e^{-E_{\Phi,0}t_\text{sep}} \nonumber \\
&+ |A_{\Phi,0}||A_{\Phi,1}|\langle 0|{\cal O}|1\rangle e^{-E_{\Phi,1}(t_\text{sep}-t)}e^{-E_{\Phi,0}t} \nonumber \\
&+ |A_{\Phi,0}||A_{\Phi,1}|\langle 1|{\cal O}|0\rangle e^{-E_{\Phi,0}(t_\text{sep}-t)}e^{-E_{\Phi,1}t} \nonumber \\
&+ |A_{\Phi,1}|^2\langle 1|{\cal O}|1\rangle e^{-E_{\Phi,1}t_\text{sep}}+ ...,
\end{align}
%
where $t_{\text{sep}}$ is the source-sink time separation, and $t$ is the gluon-operator insertion time.
The amplitudes $A_{\Phi,0}$, $A_{\Phi,1}$, and energies $E_{\Phi,0}$, $E_{\Phi,1}$ are extracted from two-state fits to the two-point correlators.
$\langle 0|{\cal O}|0\rangle$, $\langle 0|{\cal O}|1\rangle$ ($\langle 1|{\cal O}|0\rangle$), and $\langle 1|{\cal O}|1\rangle$ denote the ground-state matrix element, the ground-excited matrix element, and the excited-state matrix element, respectively.




The ground-state matrix element was obtained from a two-state fit to the three-point correlators or from a two-state simultaneous (two-sim) fit
across multiple time separations.
In order to corroborate that the fitted matrix elements were extracted correctly, the ratio
%
\begin{equation}\label{eq:Ratio}
R^\text{ratio}(z,P_z;t_\text{sep},t) = \frac{C^\text{3pt}(z,P_z;t_\text{sep},t)}{C^\text{2pt}(P_z;t_\text{sep})};
\end{equation}
%
of the three-point to the two-point correlator was computed.
The first column of Figs.~\ref{fig:a12m310_Ratio_figure} and \ref{fig:a15m310_Ratio_figure} show the kaon $R^\text{ratio}$ for the ensembles a12m310 and a15m310 for selected momenta $P_z$ and Wilson-line length $z$.
If there were no excited states, Eq.~\ref{eq:Ratio} would yield the ground-state matrix element.
It can be seen that $R^\text{ratio}$ increases as the source-sink separation increases.
In some cases, $R^\text{ratio}$ starts to converge at large separation, which shows that the contribution from excited states diminishes.
The gray bands depict the ground-state matrix elements obtained from a ``two-sim'' fit to the three-point correlators at $t_\text{sep}\in [5,9]$.
In the second and third columns, one can see the variation of the fitted ground-state matrix elements for different choices of $t_\text{sep}^\text{min}$ and $t_\text{sep}^\text{max}$.
In the examples shown in Fig.~\ref{fig:a15m310_Ratio_figure}, our ground-state fitted matrix elements are stable for the a15m310 ensemble regardless of the choices we made for $t_\text{sep}^{\text{min},\text{max}}$.
However, for the example from the a12m310 ensemble, we found that the ground-state matrix elements are sensitive to the selection of $t_\text{sep}$.
Our choices of $t_\text{sep}\in [5,9]$ for the analysis do provide a reliable ground-state extraction.



\begin{figure*}[!tbp]
  \centering
\includegraphics[width=0.35\textwidth]{images/a12m310_ratio_fit_kaon_p1z4.pdf}
\includegraphics[width=0.2\textwidth]{images/a12m310_twosimfit_kaon_p1z4_tsmin.pdf}
\includegraphics[width=0.2\textwidth]{images/a12m310_twosimfit_kaon_p1z4_tsax.pdf}\\
\includegraphics[width=0.35\textwidth]{images/a12m310_ratio_fit_kaon_p2z1.pdf}
\includegraphics[width=0.2\textwidth]{images/a12m310_twosimfit_kaon_p2z1_tsmin.pdf}
\includegraphics[width=0.2\textwidth]{images/a12m310_twosimfit_kaon_p2z1_tsmax.pdf}
\caption{Example ratio plots  (left column),
and two-sim fits (last 2 columns) from the pion mass $a\approx 0.12$~fm, $M_\pi\approx 310$~MeV for $P_z=8\pi/L$, $z=4$ (upper row) and $P_z=4\pi/L$, $z=1$ (lower row).
The gray band shown in all plots corresponds to the extracted ground-state matrix element from the two-sim fit using $t_\text{sep}\in[5,9]$.
From left to right, the columns represent:
the ratio of the three-point to two-point correlators with the reconstructed fit bands from the two-sim fit using $t_\text{sep}\in [5,9]$, shown as functions of $t-t_\text{sep}/2$,
the two-sim fit results using $t_\text{sep}\in[t_\text{sep}^\text{min},9]$ as functions of $t_\text{sep}^\text{min}$, and
the two-sim fit results using $t_\text{sep}\in[5,t_\text{sep}^\text{max}]$
as functions of $t_\text{sep}^\text{max}$.
}
\label{fig:a12m310_Ratio_figure}
\end{figure*}




\begin{figure*}[!tbp]
  \centering
\includegraphics[width=0.35\textwidth]{images/a15m310_ratio_fit_kaon_p1z4.pdf}
\includegraphics[width=0.2\textwidth]{images/a15m310_twosimfit_kaon_p1z4_tsmin.pdf}
\includegraphics[width=0.2\textwidth]{images/a15m310_twosimfit_kaon_p1z4_tsax.pdf}\\
\includegraphics[width=0.35\textwidth]{images/a15m310_ratio_fit_kaon_p2z1.pdf}
\includegraphics[width=0.2\textwidth]{images/a15m310_twosimfit_kaon_p2z1_tsmin.pdf}
\includegraphics[width=0.2\textwidth]{images/a15m310_twosimfit_kaon_p2z1_tsmax.pdf}
\caption{Example ratio plots (left column),
and two-sim fits (last 2 columns) from the pion mass $a\approx 0.15$~fm, $M_\pi\approx 310$~MeV for $P_z=1\times 2\pi/L$, $z=1$ (upper row) and $P_z=4\times 2\pi/L$, $z=4$ (lower row).
The gray band shown in all plots corresponds to the extracted ground-state matrix element from the two-sim fit using $t_\text{sep}\in[5,9]$.
From left to right, the columns represent:
the ratio of the three-point to two-point correlators with the reconstructed fit bands from the two-sim fit using $t_\text{sep}\in [5,9]$, shown as functions of $t-t_\text{sep}/2$,
the two-sim fit results using $t_\text{sep}\in[t_\text{sep}^\text{min},9]$ as functions of $t_\text{sep}^\text{min}$, and
the two-sim fit results using $t_\text{sep}\in[5,t_\text{sep}^\text{max}]$
as functions of $t_\text{sep}^\text{max}$.
}
\label{fig:a15m310_Ratio_figure}
\end{figure*}
